% Generated draft: requires mathematical and source review.
\documentclass[11pt]{article}
\usepackage[margin=25mm]{geometry}
\usepackage{fontspec}
\setmainfont{Latin Modern Roman}
\usepackage{amsmath,amssymb,mathtools,unicode-math}
\setmathfont{Latin Modern Math}
\usepackage{enumitem,needspace,xurl,hyperref}
\setlength{\parindent}{0pt}
\setlength{\parskip}{5pt}
\setlength{\emergencystretch}{3em}
\newcommand{\N}{\mathbb{N}}
\newcommand{\Z}{\mathbb{Z}}
\newcommand{\Q}{\mathbb{Q}}
\newcommand{\R}{\mathbb{R}}
\newcommand{\C}{\mathbb{C}}
\newcommand{\F}{\mathbb{F}}
\newcommand{\A}{\mathbb{A}}
\newcommand{\PP}{\mathbb{P}}
\newcommand{\E}{\mathbb{E}}
\newcommand{\eps}{\varepsilon}
\newcommand{\dd}{\,\mathrm{d}}
\begin{document}
\section*{UnsolvedMath ID: 30000604}
\textbf{Source problem code:} catalogue entry
\subsection{Definitions and mathematical statement}
Let $k,d,n$ be integers with $3\leq d\leq k$ and $(d-1)n\geq dk$, equivalently $n\geq\lceil dk/(d-1)\rceil$. Write $[n]=\{1,\ldots,n\}$ and $\binom{[n]}k$ for its $k$-element subsets. A family $\mathcal G\subseteq\binom{[n]}k$ is a $K(d)$-family if
\[
 \forall A_1,\ldots,A_d\in\mathcal G\text{ pairwise distinct}:\quad
 \left|\bigcup_{i=1}^d A_i\right|\leq2k
 \ \Longrightarrow\ 
 \bigcap_{i=1}^d A_i\ne\varnothing.
\]
This condition concerns exactly $d$ distinct members with small union; it is not silently replaced by pairwise intersection or by $d$-wise intersection with no union restriction. Equality at union size $2k$ is included.

The proposed extremal statement, throughout the stated integer range, is
\[
 |\mathcal G|\leq\binom{n-1}{k-1},
\]
with equality only when there is $x\in[n]$ such that
\[
 \mathcal G=\{A\in\tbinom{[n]}k:x\in A\}.
\]
This is the full star at a common point. Such a star satisfies the condition and has the proposed cardinality. If equality is phrased only as having a common point, its cardinality already forces the full star; no maximality hypothesis is added. No strict inequality in the threshold, large-$n$ qualification, or unspecified exceptional cases are inserted.

The supplied literature triage checked 21 August 2026 reports no authoritative complete resolution of the extracted claim in its targeted search. It supplies no theorem-author attribution, case list or endpoint refinement; the dated open label is retained without independent current-status certification.
\subsection{Short English statement}
A $k$-uniform family in which any $d$ distinct sets with union of size at most $2k$ intersect should have at most one full star's size for $n\geq dk/(d-1)$. At equality it should be exactly a full star. Both the small-union condition and the complete threshold range are preserved.
\subsection{Sources}
\begin{itemize}
\item Local UnsolvedMath registry, permanent ID 30000604, exact source code catalogue entry; \emph{Extremal Bounds for Intersecting $K(d)$-Families}. Proposed year 2006, status open; corrected\_verified statement retains $k\geq d\geq3$, the $2k$ union threshold, size bound and common-point equality case.
\item Supplied report citation: \emph{Combinatorics, Probability and Computing (2007)}, Oberwolfach Reports, DOI 10.4171/owr/2006/48, pp. 2875--2952, \url{https://doi.org/10.4171/owr/2006/48}. Title year differs from DOI/proposed year.
\item Supplied literature triage checked 2026-08-21 cites the same report, \url{https://doi.org/10.4171/owr/2006/48}, and reports no authoritative complete resolution found. The dated OWR-derived curation requests independent verification of time-sensitive status and cited claims; no separate theorem/case reference is supplied.
\end{itemize}
\end{document}
